\uuid{Xb5N}
\titre{La cellule LSTM à la main}
\niveau{L3}
\module{Apprentissage automatique}
\chapitre{Réseaux de neurones récurrents}
\sousChapitre{LSTM}
\theme{Portes d'oubli et d'entrée, cell state, mémoire à long terme}
\auteur{Maxime Nguyen}
\datecreate{2026-05-13}
\organisation{AMSCC}
\difficulte{2}

\contenu{
\texte{
  Une cellule LSTM maintient deux vecteurs d'état : la \textit{cell state} $c_t \in \mathbb{R}$ et l'état caché $h_t \in \mathbb{R}$ (on reste scalaire pour les calculs). À chaque pas de temps, quatre quantités intermédiaires sont calculées à partir de $x_t$ et $h_{t-1}$ :
  \begin{align*}
    f_t &= \sigma(w_f x_t + u_f h_{t-1} + b_f) & &\text{(porte d'oubli)}\\
    i_t &= \sigma(w_i x_t + u_i h_{t-1} + b_i) & &\text{(porte d'entrée)}\\
    \tilde{c}_t &= \tanh(w_c x_t + u_c h_{t-1} + b_c) & &\text{(candidat)}\\
    o_t &= \sigma(w_o x_t + u_o h_{t-1} + b_o) & &\text{(porte de sortie)}
  \end{align*}
  La mise à jour des états suit :
  \[
    c_t = f_t \cdot c_{t-1} + i_t \cdot \tilde{c}_t, \qquad h_t = o_t \cdot \tanh(c_t)
  \]
  où $\sigma$ désigne la fonction sigmoïde.

  \medskip
  Les paramètres sont tous les biais nuls et :

  \begin{center}
  \begin{tabular}{ccc}
    \hline
    & $w$ & $u$ \\
    \hline
    Porte d'oubli $f$ & $0$ & $0$ \\
    Porte d'entrée $i$ & $2$ & $0$ \\
    Candidat $\tilde{c}$ & $2$ & $0$ \\
    Porte de sortie $o$ & $1$ & $0$ \\
    \hline
  \end{tabular}
  \end{center}

  Les états initiaux sont $h_0 = 0$, $c_0 = 0$. On traite la séquence $x_1 = 1$, $x_2 = 1$, $x_3 = 0$, $x_4 = -1$.

  \medskip
  Valeurs utiles : $\sigma(0) = 0.5$, $\sigma(2) \approx 0.88$, $\sigma(-2) \approx 0.12$, $\tanh(2) \approx 0.96$, $\tanh(0) = 0$.
}
\begin{enumerate}
\item
  \question{(\textit{Pas $t=1$}) Calculer $f_1, i_1, \tilde{c}_1, o_1$, puis $c_1$ et $h_1$.}
  \reponse{
    \[
      f_1 = \sigma(0 \cdot 1 + 0 \cdot 0) = \sigma(0) = 0.5
    \]
    \[
      i_1 = \sigma(2 \cdot 1) = \sigma(2) \approx 0.88
    \]
    \[
      \tilde{c}_1 = \tanh(2 \cdot 1) = \tanh(2) \approx 0.96
    \]
    \[
      o_1 = \sigma(1 \cdot 1) = \sigma(1) \approx 0.731
    \]
    \[
      c_1 = 0.5 \times 0 + 0.88 \times 0.96 \approx 0.845
    \]
    \[
      h_1 = 0.731 \times \tanh(0.845) \approx 0.731 \times 0.687 \approx 0.502
    \]
  }

\item
  \question{(\textit{Pas $t=2$}) Calculer $f_2, i_2, \tilde{c}_2, o_2$, puis $c_2$ et $h_2$. Que remarque-t-on sur $c_2$ par rapport à $c_1$ ?}
  \reponse{
    Les poids $u$ sont tous nuls, donc $h_1$ n'intervient pas :
    \[
      f_2 = \sigma(0) = 0.5, \quad i_2 \approx 0.88, \quad \tilde{c}_2 \approx 0.96, \quad o_2 \approx 0.731
    \]
    \[
      c_2 = 0.5 \times 0.845 + 0.88 \times 0.96 \approx 0.423 + 0.845 \approx 1.268
    \]
    \[
      h_2 = 0.731 \times \tanh(1.268) \approx 0.731 \times 0.852 \approx 0.623
    \]
    La cell state $c_2 > c_1$ : la porte d'entrée a de nouveau injecté de l'information positive, qui s'accumule dans $c_t$ grâce à la porte d'oubli (qui conserve environ la moitié de l'état précédent à chaque pas).
  }

\item
  \question{(\textit{Pas $t=3$}) L'entrée est $x_3 = 0$. Calculer $c_3$ et $h_3$. Que fait la porte d'oubli ici ? Comparer $c_3$ et $c_2$.}
  \reponse{
    \[
      f_3 = \sigma(0) = 0.5, \quad i_3 = \sigma(0) = 0.5, \quad \tilde{c}_3 = \tanh(0) = 0, \quad o_3 = \sigma(0) = 0.5
    \]
    \[
      c_3 = 0.5 \times 1.268 + 0.5 \times 0 = 0.634
    \]
    \[
      h_3 = 0.5 \times \tanh(0.634) \approx 0.5 \times 0.560 \approx 0.280
    \]
    Avec une entrée nulle, la porte d'oubli conserve la moitié de $c_2$ et la porte d'entrée n'ajoute rien ($\tilde{c}_3 = 0$). La cell state diminue donc de moitié : $c_3 \approx c_2 / 2$. Le réseau « oublie » progressivement l'information accumulée.
  }

\item
  \question{(\textit{Pas $t=4$}) L'entrée est $x_4 = -1$. Calculer $c_4$ et $h_4$. Commenter le rôle de la porte d'entrée.}
  \reponse{
    \[
      f_4 = \sigma(0) = 0.5, \quad i_4 = \sigma(-2) \approx 0.12, \quad \tilde{c}_4 = \tanh(-2) \approx -0.96, \quad o_4 = \sigma(-1) \approx 0.269
    \]
    \[
      c_4 = 0.5 \times 0.634 + 0.12 \times (-0.96) \approx 0.317 - 0.115 \approx 0.202
    \]
    \[
      h_4 = 0.269 \times \tanh(0.202) \approx 0.269 \times 0.200 \approx 0.054
    \]
    La porte d'entrée $i_4 \approx 0.12$ est faible pour une entrée négative : le candidat négatif $\tilde{c}_4$ n'est injecté que partiellement dans $c_t$. Cela atténue l'effondrement brutal de la cell state, qui reste positive. La porte d'entrée régule ainsi la quantité d'information nouvelle incorporée selon le signe de l'entrée.
  }

\item
  \question{Dans l'exercice 1, l'état caché $h_t$ était directement écrasé à chaque pas. Expliquer en quoi le mécanisme de cell state $c_t$ de la LSTM permet de mieux préserver l'information sur le long terme, en lien avec l'exercice 3.}
  \reponse{
    Dans le RNN simple, $h_t = \tanh(\cdots)$ est une transformation non linéaire saturante de $h_{t-1}$ : l'état précédent est systématiquement « compressé » dans $[-1, 1]$, ce qui fait que $\frac{\partial h_t}{\partial h_{t-1}} = w_h(1 - h_t^2)$ peut être proche de $0$ si $h_t$ est saturé (exercice 3).

    Dans la LSTM, la mise à jour de $c_t$ est \emph{additive} et modulée par des portes :
    \[
      c_t = f_t \cdot c_{t-1} + i_t \cdot \tilde{c}_t
    \]
    Si $f_t \approx 1$ et $i_t \approx 0$, alors $c_t \approx c_{t-1}$ : la cell state est copiée presque intacte, et $\frac{\partial c_t}{\partial c_{t-1}} \approx 1$. Le gradient peut ainsi se propager sans s'atténuer sur de longues séquences, ce qui atténue le problème du \textit{gradient vanishing}.
  }
\end{enumerate}
}
